\section{Билинейные формы}

\begin{definition}
    Пусть $V$ --- линейное пространство над полем $k$.
    \emph{Билинейная форма} --- это билинейное отображение $B\colon V \times V \to k$.
\end{definition}

\begin{remark}
    Многие дальнейшие утверждения требуют $\operatorname{char}(k) \ne 2$, так что будем считать это условие выполненным по умолчанию.
\end{remark}


\begin{definition}
    По билинейной форме $B$ определяются две \emph{корреляции}:

    \emph{Правая корреляция:}
    \[
        B^{\wedge}\colon V \to V^*, \quad v \mapsto (u \mapsto B(u, v)).
    \]

    \emph{Левая корреляция:}
    \[
        {}^{\wedge}\!B\colon V \to V^*, \quad v \mapsto (u \mapsto B(v, u)).
    \]

    Обе корреляции --- линейные отображения.
\end{definition}


\begin{definition}
    \emph{Правый ортогонал} (правое ядро):
    \[
        V^\perp = \ker B^{\wedge} = \{v \in V \mid B(u, v) = 0\ \forall u \in V\}.
    \]
    \emph{Левый ортогонал} (левое ядро):
    \[
        {}^\perp V = \ker {}^{\wedge}\!B = \{v \in V \mid B(v, u) = 0\ \forall u \in V\}.
    \]
\end{definition}


\begin{remark}
    Если $B$ симметрична или кососимметрична, то $V^\perp = {}^\perp V$, и можно говорить просто о ядре $\ker B$.
\end{remark}


\begin{definition}[Матрица Грама]
    Пусть $\mathbf{v} = (v_1, \dots, v_n)$ и $\mathbf{u} = (u_1, \dots, u_m)$ --- два набора векторов в
    $V$. Тогда
    \[
        G_{\mathbf{v}, \mathbf{u}} = \bigl(B(v_i, u_j)\bigr) \in \mathrm{Mat}_{n \times m}(k).
    \]
    Таким образом
    \[
        B(x, y) = B\!\left(\sum_i x_i e_i,\ \sum_j y_j e_j\right)
        = \sum_{i,j} x_i y_j B(e_i, e_j)
        = \sum_{i,j} x_i G_{ij} y_j
        = x^T G y.
    \]
\end{definition}

\begin{definition}[Изометрия]
    Пусть $(V, B)$ и $(W, \widetilde{B})$ --- пространства с билинейными формами. Линейное отображение
    $f\colon V \to W$ называется \emph{изометрией}, если
    \[
        B(v_1, v_2) = \widetilde{B}(f(v_1), f(v_2)) \quad \forall\, v_1, v_2 \in V.
    \]
    Такое отображение -- изометрический  изоморфизм, если является биекцией.
\end{definition}


\begin{stmt}
    Если $f\colon V \to W$ --- биекция и изометрия, то $f^{-1}$ тоже изометрия.
    \begin{proof}
        $\widetilde{B}(w_1, w_2) = \widetilde{B}(f(f^{-1}(w_1)),\, f(f^{-1}(w_2))) = B(f^{-1}(w_1),\, f^{-1}(w_2))$.
    \end{proof}
\end{stmt}
